% Part: many-valued-logic
% Chapter: syntax-and-semantics
% Section: semantic-notions

\documentclass[../../../include/open-logic-section]{subfiles}

\begin{document}

\olfileid{mvl}{syn}{sem}

\olsection{مفاهيم دلالية}

إذا أعطينا منطقًا متعدد القيم~$\Log L$ بمصفوفته، أمكن أن ننقل إليه
المفاهيم الدلالية المألوفة لـ~$\Log L$ على الوجه الآتي.

\begin{defn}
\begin{enumerate}
\item تكون~!!^a{formula}~$!A$ \emph{قابلة للإشباع} إذا وجد
  $\pAssign{\OLId{latin-ordinary}{v}}$ يحقق~$\pSat{\OLId{latin-ordinary}{v}}{!A}$؛ وتكون \emph{غير قابلة
  للإشباع} إذا لم يوجد~$\pAssign{\OLId{latin-ordinary}{v}}$ يحقق~$\pSat{\OLId{latin-ordinary}{v}}{!A}$؛
\item وتكون~!!^a{formula}~$!A$ \emph{تحصيلًا منطقيًا} إذا تحقق
  $\pSat{\OLId{latin-ordinary}{v}}{!A}$ لجميع~!!{valuation}s~$\OLId{latin-ordinary}{v}$؛
\item وإذا كانت~$\OLId{greek-variable}{Gamma}$ مجموعة~!!{formula}s، كان
  $\OLId{greek-variable}{Gamma} \Entails !A$، أي «$\OLId{greek-variable}{Gamma}$ تستلزم~$!A$»، إذا وفقط إذا
  تحقق~$\pSat{\OLId{latin-ordinary}{v}}{!A}$ لكل~!!{valuation}~$\pAssign{\OLId{latin-ordinary}{v}}$ يحقق
  $\pSat{\OLId{latin-ordinary}{v}}{\OLId{greek-variable}{Gamma}}$.
\item وإذا كانت~$\OLId{greek-variable}{Gamma}$ مجموعة~!!{formula}s، كانت~$\OLId{greek-variable}{Gamma}$
  \emph{قابلة للإشباع} إذا وجد~!!a{valuation}~$\pAssign{\OLId{latin-ordinary}{v}}$ يحقق
  $\pSat{\OLId{latin-ordinary}{v}}{\OLId{greek-variable}{Gamma}}$؛ وإلا
  كانت~$\OLId{greek-variable}{Gamma}$ \emph{غير قابلة للإشباع}.
\end{enumerate}
\end{defn}

ولهذه المفاهيم جملة من الخواص التي عرفناها في المنطق الكلاسيكي:

\begin{prop}
\ollabel{prop:semanticalfacts}
\begin{enumerate}
\item تكون~$!A$ تحصيلًا منطقيًا إذا وفقط إذا كان
  $\emptyset \Entails !A$؛
\item إذا كانت~$\OLId{greek-variable}{Gamma}$ قابلة للإشباع، كانت كل مجموعة جزئية منتهية
  من~$\OLId{greek-variable}{Gamma}$ قابلة للإشباع أيضًا؛
\item\ollabel{def:monotonicity}%
الرتابة: إذا كان~$\OLId{greek-variable}{Gamma} \subseteq \OLId{greek-variable}{Delta}$ و$\OLId{greek-variable}{Gamma} \Entails !A$،
  كان~$\OLId{greek-variable}{Delta} \Entails !A$ أيضًا؛
\item\ollabel{def:Cut}%
التعدي: إذا كان~$\OLId{greek-variable}{Gamma} \Entails !A$ و
  $\OLId{greek-variable}{Delta} \cup \{ !A\} \Entails !B$، كان
  $\OLId{greek-variable}{Gamma} \cup \OLId{greek-variable}{Delta} \Entails !B$؛
\end{enumerate}
\end{prop}

\begin{proof}
يترك إثباتها تمرينًا.
\end{proof}

\begin{prob}
أثبت \olref[mvl][syn][sem]{prop:semanticalfacts}.
\end{prob}

في المنطق الكلاسيكي صلات وثيقة بين الاستلزام والشرط. فمنها صحة
\emph{قاعدة الفصل}: إذا كان~$\OLId{greek-variable}{Gamma} \Entails !A$ و
$\OLId{greek-variable}{Gamma} \Entails !A \lif !B$، لزم~$\OLId{greek-variable}{Gamma} \Entails !B$.
ومنها مبرهنة الاستنباط الدلالية التي تصل~$\Entails$ بـ~$\lif$:
$\OLId{greek-variable}{Gamma} \Entails !A \lif !B$ إذا وفقط إذا كان
$\OLId{greek-variable}{Gamma} \cup \{!A\} \Entails !B$. لكن هذه النتائج \emph{لا} تلزم
في كل منطق متعدد القيم؛ إذ يتوقف صدقها على دالة الصدق~$\tf{\lif}$.

\end{document}
